\documentclass[11pt,a4paper]{article}

\usepackage{fontspec}
\usepackage{xeCJK}
\setCJKmainfont{Hiragino Mincho ProN}
\setCJKsansfont{Hiragino Kaku Gothic ProN}
\usepackage[margin=2.5cm]{geometry}
\usepackage{amsmath,amssymb,amsfonts}
\usepackage{graphicx}
\usepackage{hyperref}
\usepackage{booktabs}
\usepackage{amsthm}

\newcommand{\Spec}{\mathrm{Spec}}
\newcommand{\Zp}{\mathbb{Z}[1/p]}

\begin{document}

\title{素数チャンネルミュートの物理的実装：\\算術的真空工学のための10の実験プラットフォーム}
\author{Wright Brothers\\{\small 独立研究}}
\date{2026年4月}
\maketitle

\begin{abstract}
関連論文~\cite{paper_A}は、素数$p$の倍数にあたる真空モードを抑制する操作（「素数チャンネルミュート」）がゼータ正則化真空エネルギーの符号を反転させることを示し、単一の実験プラットフォーム（SQUIDノッチフィルタを備えた超伝導回路）を提案した。本論文では、素数チャンネルミュートを実現するための\emph{全て}の実行可能な物理プラットフォームを、ボスト-コンヌ系の枠組みで特定された3つのメカニズム —— (A)~スペクトルフィルタリング、(B)~算術的境界条件、(C)~相転移セクター選択 —— に沿って体系的に調査する。10の異なる実装候補を特定し、5つの基準（技術成熟度、対応可能素数数、真空エネルギー測定の直接性、コスト、タイムライン）で評価する。さらに、最小限のリソースで即座に実行可能な3つの実験を提案する：音響共振器による原理実証（1万円、1日）、クラウド量子ハードウェア上でのボスト-コンヌ・ハミルトニアンの量子シミュレーション（0円、1週間）、およびFabry-P\'{e}rotフィルタリングを備えた光学周波数コム実験（100万円、6ヶ月）。これらは~\cite{paper_A}のSQUID実験を補完し、異なる物理領域にわたる真空エネルギーの算術構造の独立した検証を可能にする。
\end{abstract}

% ============================================================================
\section{序論}
\label{sec:intro}
% ============================================================================

\cite{paper_A}において、任意の素数$p$をミュートするとゼータ正則化真空エネルギーの符号が反転すること——$\zeta_{\neg p}(-3) = (1-p^3)/120 < 0$——を証明した。射影$P_{\neg p} = \sum_{p\nmid n}|n\rangle\langle n|$を実現する3つの物理メカニズムを特定した：
\begin{itemize}
\item[(A)] スペクトルフィルタリング：$p$の倍数周波数のモードを吸収
\item[(B)] 算術的境界条件：$\gcd(n,p)=1$のモードのみを許容する境界条件
\item[(C)] 相転移セクター選択：BC系の$\mathbb{Z}_p^*$不変KMS状態を選択
\end{itemize}

\cite{paper_A}ではメカニズム(A)をSQUIDベースの詳細な実験に展開した。本論文では全メカニズムを体系的に調査し、9つの追加プラットフォームを特定する。うち3つは最小限のリソースで即座に実行可能である。

% ============================================================================
\section{メカニズムA：スペクトルフィルタリング}
\label{sec:mech_a}
% ============================================================================

メカニズムAは、$p$の倍数にあたる周波数成分を吸収して対応するモードを抑制する。物理系のスペクトルの種類に応じて2つのサブケースを区別する。

\subsection{A1：対数スペクトル系}

BC系では$E_n = \log(n)$である。$\omega = \log(p)$の単一のノッチフィルタが全ての$p$可除状態を除去する。これは$p|n$なる任意の$n$が素因数分解の寄与に$\log(p)$を含むためである。

\textbf{プラットフォーム1：対数ポテンシャル荷電ワイヤー。}
2D荷電ワイヤーはポテンシャル$V(r) = V_0\log(r/a)$を生成する。WKB近似で束縛状態のエネルギーは$E_n \sim \log(n)$とスケールする。$\omega = \log(p)$の周波数選択的吸収体を備えたこのポテンシャル中の粒子は、素数ミュートを伴うBCハミルトニアンを実現する。

技術成熟度：低。タイムライン：3--5年。

\subsection{A2：等間隔スペクトル系}

標準的な空洞は$\omega_n = n\omega_0$（等間隔）を持つ。素数$p$をミュートするには、$p\omega_0$, $2p\omega_0$, $3p\omega_0$, …のモードを抑制する必要がある——これは周期$p\omega_0$の\emph{周期コムフィルタ}である。

\textbf{プラットフォーム2：超伝導回路 + SQUIDフィルタ}（\cite{paper_A}に詳述）。同一平面導波路共振器（$f_0 \approx 6$\,GHz）にSQUIDベースのノッチフィルタを各偶数高調波に同調させる。技術成熟度：高。コスト：500万〜1000万円。タイムライン：12ヶ月。

\textbf{プラットフォーム3：光学周波数コム + Fabry-P\'{e}rot。}
モードロックレーザーが$f_n = f_{\mathrm{CEO}} + n f_{\mathrm{rep}}$（$f_{\mathrm{rep}} \sim 100$\,MHz）の周波数コムを生成する。自由スペクトル幅$\mathrm{FSR} = p \cdot f_{\mathrm{rep}}$のFabry-P\'{e}rot共振器は$p$の倍数のモードのみを透過し、それ以外を反射する。

\emph{反射}ビームは$p$と互いに素なモード$\{n : p\nmid n\}$のみを含む。これが$p$-ミュートされたスペクトルである。

測定：反射ビームの量子ノイズ（ショットノイズ）がフィルタされたモード集合の真空揺らぎをエンコードする。フィルタあり/なしのノイズパワーの比較により、真空揺らぎエネルギーがモード選択の算術に依存するか検証する。

技術成熟度：高（周波数コムはノーベル賞技術~\cite{Hansch2005}；Fabry-P\'{e}rot共振器は標準的）。コスト：100万〜500万円。タイムライン：6ヶ月。

\textbf{プラットフォーム4：音響共振管 + 吸音材。}
長さ$L$の管は$f_n = nv/(2L)$の定在波を支持する。位置$x = L/(2p)$に吸音材を挿入すると、その位置に腹を持つモードが優先的に減衰する。

$p = 2$の場合：$L/4$に吸音材を置くと奇数モード（$L/4$で腹）が減衰し、偶数モード（$L/4$で節）が生存する。これは$p$ミュートの\emph{補集合}（偶数が残り奇数が消える）だが、原理は同一であり、相補的な測定として同等に有益である。

技術成熟度：自明（アクリル管、スピーカー、マイクロフォン）。コスト：約1万円。タイムライン：\textbf{1日}。算術的モード選別の最低コスト原理実証。

\textbf{プラットフォーム5：ピエゾ圧電素子アレイ。}
ピエゾ結晶（水晶、PZT）は鋭い機械的共鳴を持つ。共振周波数$f_n = nf_0$の$N$個の結晶アレイに$p$の倍数での選択的減衰を施すことで、音響素数ミュートの固体プラットフォームを提供する。

技術成熟度：中程度。コスト：約50万円。タイムライン：3ヶ月。

% ============================================================================
\section{メカニズムB：算術的境界条件}
\label{sec:mech_b}
% ============================================================================

メカニズムBは、カシミール効果で導体板が特定波長のモードを選択するのと同様に、$p$と互いに素なモードを幾何学的に選択する境界条件を課す。

\textbf{プラットフォーム6：半導体超格子。}
GaAs/AlGaAs超格子を周期$d = p\cdot a_0$（$a_0$は基本格子定数）で形成すると、$p$の倍数のモード番号にミニバンドギャップが開く。許容ミニバンド中の電子は$\gcd(n,p) = 1$のモードのみを占有する。

$p = 2$の場合：$a_0 = 5$\,nm, $d = 10$\,nm, ギャップエネルギー$\sim 0.1$\,eV。偶数番号のミニバンドにギャップ；奇数番号は伝搬。

技術成熟度：高（MBE成長は確立技術）。コスト：約500万円。タイムライン：1年。

\textbf{プラットフォーム7：電磁メタマテリアル。}
周期$p \cdot a$のスプリットリング共振器を持つ3Dメタマテリアルは、$p$の倍数で電磁バンドギャップを開く。3つの独立な格子方向で3つの独立な素数に対応可能：$x$方向に周期$2a$、$y$方向に$3a$、$z$方向に$5a$で$p = 2, 3, 5$を同時ミュート。

技術成熟度：中程度。コスト：約2000万円。タイムライン：2年。

\textbf{プラットフォーム8：オイラー積結晶。}
メカニズムBの最も野心的なバージョン：素数$p_1, p_2, p_3, \ldots$に対してそれぞれ周期$p_i \cdot a$の複数のブラッグ格子を積層する。

層1（周期$2a$）：2の倍数にギャップ。
層2（周期$3a$）：3の倍数にギャップ。
層3（周期$5a$）：5の倍数にギャップ。

この組み合わせは\emph{エラトステネスの篩}の物理的実現：素数番号のモード（$n=1$を含む）のみが全層を通過する。有効ゼータ関数は$\zeta(s)\prod_i(1-p_i^{-s})$となり、複数のオイラー因子が同時に除去される。

数値見積もり：最初の3素数（$2, 3, 5$）で$n \leq 100$のモードの77\%が除去される。

技術成熟度：低。コスト：約5000万円。タイムライン：3--5年。

% ============================================================================
\section{メカニズムC：相転移セクター選択}
\label{sec:mech_c}
% ============================================================================

メカニズムCはBC相転移を直接操作し、$\mathbb{Z}_p^*$不変KMS状態を選択することで素数ミュートを実装する。

\textbf{プラットフォーム9：量子コンピュータシミュレーション。}
BCハミルトニアン$H|n\rangle = \log(n)|n\rangle$を$\ell^2(\{1,\ldots,2^N\})$上に$N$量子ビットで実装する。素数ミュート摂動$H' = H + \Delta\sum_{p|n}|n\rangle\langle n|$は$\Delta \to \infty$の極限で$p$可除状態をギャップアウトさせる。

プロトコル：(1) VQEまたは断熱準備で$H$の基底状態を準備。(2) $\langle H\rangle$を測定。(3) $\Delta$を増加させながら摂動$H'$を追加。(4) $\Delta$の関数として$\langle H'\rangle$を測定。(5) $\Delta \to \infty$で$\langle H'\rangle \to \zeta_{\neg p}(-1)$への収束を確認。

IBM QuantumまたはGoogle Cirqのクラウドプラットフォームで5--10量子ビットにより実行可能であり、最初の32--1024モードをカバーする。

技術成熟度：高。コスト：\textbf{0円}（クラウド無料枠）。タイムライン：\textbf{1週間}。

\textbf{プラットフォーム10：冷却原子光格子。}
1D光格子にサイト依存ポテンシャル$V_n = V_0\log(n)$（整形レーザービームで近似）で超冷却原子を装填する。深さ$\Delta$の周期$p$の二次格子が$p$可除サイトのエネルギーをシフトする。

$p = 2$の場合：二次格子の波長$\lambda_2 = 2\times\lambda_1 = 1064$\,nm（Nd:YAG基本波）。$p = 3$：$\lambda_3 = 1596$\,nm。$p = 5$：$\lambda_5 = 2660$\,nm。全て標準的なレーザー波長。

冷却原子プラットフォームは数千の格子点と複数素数の同時対応が可能であり、BC系の最もスケーラブルなアナログ量子シミュレータである。

密接に関連するプラットフォームとして、空間的に変調された音速を持つボーズ-アインシュタイン凝縮体（BEC）がある。BECフォノンは「アナログ重力」を実現するため~\cite{Unruh1981}、算術的変調を施したBECは\emph{算術的に修正された真空におけるアナログ重力}を実装する——代数的メカニズムから時空物理への最も直接的な橋渡しである。

技術成熟度：中程度。コスト：約5000万円（既存ラボ）〜約2億円（新規セットアップ）。タイムライン：2--3年。

% ============================================================================
\section{比較評価}
\label{sec:comparison}
% ============================================================================

表\ref{tab:comparison}に全10プラットフォームを要約する。

\begin{table*}[t]
\centering
\caption{素数チャンネルミュートのための10の実装プラットフォーム。TRL=技術成熟度（1--5段階）。$N_p$=同時対応可能素数数。測定：D=真空エネルギー直接測定、I=間接（スペクトル）。}
\label{tab:comparison}
\small
\begin{tabular}{@{}rlcccccc@{}}
\toprule
\# & プラットフォーム & 機構 & TRL & $N_p$ & 測定 & コスト & 期間 \\
\midrule
1 & 対数ポテンシャルトラップ & A1 & 2 & 1--2 & I & 1000万円 & 3--5年 \\
2 & 超伝導回路+SQUID~\cite{paper_A} & A2 & 4 & 2--5 & D & 500万〜1000万円 & 1年 \\
3 & 周波数コム+Fabry-P\'{e}rot & A2 & 4 & 1--3 & I & 100万〜500万円 & 6ヶ月 \\
4 & 音響共振管 & A2 & 5 & 1--2 & I & 1万円 & 1日 \\
5 & ピエゾ圧電素子アレイ & A2 & 4 & 1--2 & I & 50万円 & 3ヶ月 \\
6 & 半導体超格子 & B & 4 & 1--2 & I & 500万円 & 1年 \\
7 & 電磁メタマテリアル & B & 3 & 1--3 & I & 2000万円 & 2年 \\
8 & オイラー積結晶 & B & 2 & 3--5 & I & 5000万円 & 3--5年 \\
9 & 量子コンピュータ（クラウド） & C & 5 & 多数 & I & 0円 & 1週間 \\
10 & 冷却原子光格子 & C & 3 & 3--5 & D & 5000万円 & 2--3年 \\
\bottomrule
\end{tabular}
\end{table*}

% ============================================================================
\section{即座に実行可能な3つの実験}
\label{sec:immediate}
% ============================================================================

\cite{paper_A}のSQUID実験を補完する、即座に開始可能な3つの実験を詳述する。

\subsection{実験$\alpha$：音響共振管（1万円、1日）}

\textbf{材料}：アクリル管（$L = 1$\,m）、スピーカー、マイクロフォン、FFTソフトウェアを備えたPC。

\textbf{プロトコル}：(1) ホワイトノイズで管を駆動。(2) FFTで共鳴スペクトルを記録：$f_n = nv/(2L) \approx n \times 172$\,Hzのピーク群。(3) $L/2$に吸音フォームを挿入：奇数モード（$L/2$で腹）が減衰；偶数モード（$L/2$で節）が生存。(4) $L/3$に吸音フォームを挿入：3の倍数でないモードがさらに減衰。(5) 挿入前後のスペクトルピーク高さを比較。

\textbf{検証事項}：算術的位置の物理的吸音材が可除性に基づいてモードを選択的に除去すること。古典（非量子）的な算術的モード選別の実証。

\textbf{限界}：真空エネルギーは調べられない（古典音響）。しかし、物理的幾何構造が算術的フィルタリングを実装するという原理を確立する。

\subsection{実験$\beta$：量子シミュレーション（0円、1週間）}

\textbf{プラットフォーム}：IBM Quantum（無料クラウドアクセス）またはGoogle Cirq。

\textbf{プロトコル}：(1) 対角ゲートを使用してBCハミルトニアン$H|n\rangle = \log(n)|n\rangle$を5--10量子ビットにエンコード。(2) VQEまたは断熱準備で逆温度$\beta$での熱状態を準備。(3) $\langle H\rangle = -\partial_\beta \log Z(\beta)$を測定。(4) $\Delta$を増加させながら$p = 2$の摂動$\Delta\sum_{p|n}|n\rangle\langle n|$を追加。(5) $\Delta$の関数として$\langle H\rangle$を測定。(6) $\Delta \to \infty$で$\zeta_{\neg 2}(-1) = +1/12$への収束を確認（$\zeta(-1) = -1/12$からの符号反転）。

\textbf{検証事項}：素数ミュートを伴うBC系が予測された分配関数変更を再現すること。実際の量子ハードウェア上での代数的メカニズムの検証。

\textbf{限界}：BC系のシミュレーションであり、物理的真空ではない。物理的真空がBC構造を持つことの証明にはならない。

\subsection{実験$\gamma$：光学周波数コム（100万〜500万円、6ヶ月）}

\textbf{材料}：モードロックレーザー（$f_{\mathrm{rep}} \sim 100$\,MHz）、Fabry-P\'{e}rot共振器（$\mathrm{FSR} = 2f_{\mathrm{rep}}$）、バランスドホモダイン検出器。

\textbf{プロトコル}：(1) 周波数コム$f_n = f_{\mathrm{CEO}} + nf_{\mathrm{rep}}$を生成。(2) $\mathrm{FSR} = 2f_{\mathrm{rep}}$のFabry-P\'{e}rotに通過：透過ビームは偶数モードのみ；反射ビームは奇数モードのみ（$p = 2$ミュート）。(3) 反射ビームの量子ノイズパワー（ショットノイズ限界 = 真空揺らぎエネルギー）を測定。(4) 非フィルタビームのノイズと比較。(5) $\mathrm{FSR} = 3f_{\mathrm{rep}}$（$p = 3$ミュート）、$\mathrm{FSR} = 5f_{\mathrm{rep}}$（$p = 5$ミュート）で繰り返し。(6) 乗法構造を検証：ノイズ比$R(p_1, p_2) = R(p_1) \times R(p_2)$？

\textbf{検証事項}：フィルタされたモード集合の真空揺らぎノイズがオイラー積と整合する形でフィルタの算術に依存するかどうか。

\textbf{限界}：真空「エネルギー」ではなく真空「揺らぎ」を測定。ノイズパワーとゼータ正則化エネルギーの接続にはショットノイズとゼロ点エネルギーを関連付ける量子光学の形式~\cite{Caves1981}を要する。

% ============================================================================
\section{スケーリングに関する考察}
\label{sec:scaling}
% ============================================================================

\cite{paper_B}において、弱エネルギー条件は連続的なエネルギー閾値ではなく離散的な位相不変量（レギュレータの方向）であると予想した。正しければスケーリングは不要：単一の局所化$\mathbb{Z} \to \Zp$がエネルギースケールに依存せずWECを反転させる。

この予想とは独立に、各プラットフォームには自然なスケーリング経路がある：
\begin{itemize}
\item プラットフォーム2（SQUID）：$N$個の同期共振器の配列、真空エネルギー変更がコヒーレントなら$E \propto N^2$。
\item プラットフォーム6（超格子）：バルク材料は体積に比例。
\item プラットフォーム8（オイラー積結晶）：3D構造、除去されたオイラー因子数に比例するエネルギー密度。
\item プラットフォーム10（光格子/BEC）：現在の実験で$\sim 10^6$サイト；フォノン媒介相互作用が自然な長距離コヒーレンスを提供。
\end{itemize}

BECプラットフォーム（プラットフォーム10の変種）はスケーリングに特に有望である。(a) BECフォノンは確立されたアナログ重力プラットフォーム~\cite{Unruh1981}；(b) フェッシュバッハ共鳴による音速の精密チューニングが算術的周期での空間変調を可能にする；(c) BECのモード数（$\sim 10^6$）はマイクロ波回路（$\sim 20$）をはるかに超える。

% ============================================================================
\section{議論}
\label{sec:discussion}
% ============================================================================

10のプラットフォームはコストで4桁（0円〜2億円）、タイムラインで5桁（1日〜5年）にわたる。3つの全メカニズム(A, B, C)と2つのスペクトル領域（マイクロ波と光学）をカバーする。

推奨するマルチプラットフォーム実験戦略：
\begin{enumerate}
\item \textbf{即時}（1ヶ月目）：実験$\alpha$（音響、原理検証）と$\beta$（量子シミュレーション、BC系検証）を同時実行。コスト：1万円 + 0円 = 1万円。
\item \textbf{短期}（1--6ヶ月目）：実験$\gamma$（周波数コム、真空揺らぎ検証）を遂行。コスト：100万〜500万円。
\item \textbf{中期}（6--18ヶ月目）：\cite{paper_A}のSQUID実験（プラットフォーム2、真空エネルギー直接測定）を実施。コスト：500万〜1000万円。
\item \textbf{長期}（2--5年目）：半導体超格子（プラットフォーム6）または冷却原子（プラットフォーム10）で多重素数ミュートとスケーリング研究を展開。
\end{enumerate}

実験$\alpha$と$\beta$が成功し（原理とBC代数構造を確認）、実験$\gamma$が真空ノイズの算術的依存性を示せば、より高コストなプラットフォーム2, 6, 10への投資の根拠が強固になる。

% ============================================================================
\section{結論}
\label{sec:conclusion}
% ============================================================================

素数チャンネルミュートを実装するための10の物理プラットフォームを特定した。音響、電磁、超伝導、半導体、メタマテリアル、量子コンピュータ、冷却原子の技術にまたがる。3つの実験は最小コスト（それぞれ1万円、0円、100万円）で即座に開始可能であり、古典音響、デジタル量子シミュレーション、量子光学にわたる算術的真空エネルギー仮説の独立した検証を提供する。

プラットフォームの多様性は算術的真空工学プログラムの強みである：真空エネルギーのオイラー積構造が（正則化の人工物ではなく）物理法則であるならば、量子化されたモードスペクトルを持つ\emph{全ての}物理系において、物理基質によらず発現するはずである。マルチプラットフォームでの確認は最も強力な証拠を提供する。

\section*{謝辞}
本研究はWright Brothers共同研究グループにより独立研究として遂行された。

\begin{thebibliography}{20}
\bibitem{paper_A} Wright Brothers, ``Arithmetic vacuum engineering: Prime channel muting via Bost-Connes phase transitions,'' companion paper (2026).
\bibitem{paper_B} Wright Brothers, ``Sheaf-theoretic reformulation of vacuum energy conditions on $\Spec(\mathbb{Z})$,'' companion paper (2026).
\bibitem{Hansch2005} T.~W.~H\"ansch, Rev. Mod. Phys. \textbf{78}, 1297 (2006).
\bibitem{Unruh1981} W.~G.~Unruh, Phys. Rev. Lett. \textbf{46}, 1351 (1981).
\bibitem{Caves1981} C.~M.~Caves, Phys. Rev. D \textbf{23}, 1693 (1981).
\bibitem{Wilson2011} C.~M.~Wilson et al., Nature \textbf{479}, 376 (2011).
\end{thebibliography}

\end{document}
